\chapter{Дүгнэлт}

Энэхүү судалгаагаар анги танхимын орчинд ажиллах камер болон нүүр танихад суурилсан ирц бүртгэлийн системийн prototype боловсруулж, бодит өгөгдөл дээр туршин үнэлэв. Туршилтын явцад нүүр илрүүлэх, эмбеддинг үүсгэх, таних, мөн ирцийг олон ажиглалтаар баталгаажуулах шатлал нь практик орчинд хэрэгжих боломжтойг харуулсан.

\section{Судалгааны асуултуудын хариу}

Нэгдүгээрт, системийн нарийвчлалын хувьд enrollment өгөгдөл дээр leave-one-out validation хийхэд hybrid болон k-NN арга 98.35\% gated accuracy, cosine арга 95.05\% gated accuracy, SVM арга 39.38\% gated accuracy үзүүлэв. Энэ нь жижиг ангид зориулсан бодит тохируулгатай нөхцөлд hybrid болон k-NN нь хамгийн тогтвортой шийдэл болохыг харуулж байна.

Хоёрдугаарт, бодит ажиллагааны хугацааны хувьд 546 enrollment зураг дээр gallery rebuild 460.10 секундэд дууссан бөгөөд 56 зурагтай туршилтын багц дээр хийсэн MTCNN танилт 67.50 секунд үргэлжиллээ. Энэ нь дунд зэргийн хэмжээтэй ангийн өгөгдөлд системийн боловсруулалтын хугацаа бодит хэрэглээнд боломжийн түвшинд байгааг харуулна. Нэг зураг боловсруулах дундаж хугацаа 1.21 секунд, нэг нүүрний мөр боловсруулах дундаж хугацаа ойролцоогоор 259 миллисекунд байв.

Гуравдугаарт, бодит хичээлийн сешнүүд дээр системийн эцсийн ирцийн гарц жигд биш байв. 3-р паарын нэг сешнд 19 сурагч, хичээлээс гадуурх үеийн нэг сешнд 14 сурагч ирсэн гэж баталгаажсан бол 2-р болон 4-р паарын сешнүүдэд тухайн тохиргоогоор ирээгүй гэж буцаасан. Энэ нь гэрэлтүүлгийн нөхцөл, sampling-ийн нягтрал, босго тохируулга гурвын нөлөө хүчтэй байгааг харуулж байна.

Дөрөвдүгээрт, практик хязгаарлалт тодорхой болов. 22 сурагчтай нэг ангийн өгөгдөл нь судалгааны анхны нотолгоо өгөхөд хангалттай боловч эцсийн нийтлэг дүгнэлт гаргахад бага хэмжээтэй. Мөн ground truth manual attendance CSV бүрэн архивлагдаагүй тул FAR, FNR, болон статистикийн ач холбогдлын бүрэн баталгаажуулалт энэ хувилбарт хязгаарлагдмал хэвээр байна.

Иймээс enrollment leave-one-out дээрх 98.35\% gated accuracy-г бодит ангийн attendance accuracy гэж шууд тайлбарлахгүй байх нь судалгааны үнэн зөвийн үндсэн шаардлага юм. Эхнийх нь таних pipeline-ийн дотоод ялгах чадварыг, хоёр дахь нь бодит орчны эцсийн ирц баталгаажуулалтын чанарыг хэмждэг өөр түвшний үзүүлэлтүүд болно.

\section{Ерөнхий дүгнэлт}

Судалгааны үндсэн зорилго болох камер ашиглан царай танихад суурилсан ирц бүртгэлийн загвар боловсруулах, турших, үнэлэх ажил биелсэн. Hybrid болон k-NN нь энэхүү өгөгдлийн хүрээнд хамгийн сайн гүйцэтгэл үзүүлсэн бөгөөд MTCNN-ийг үнэлгээний урсгалд ашиглах нь RetinaFace-тай харьцуулахад хурд болон практик хэрэгжилтийн хувьд давуу байсныг туршилт харуулсан.

Өмнөх судалгаануудтай харьцуулахад, FaceNet болон ArcFace төрлийн ажилууд LFW зэрэг хяналттай benchmark дээр 99\%-аас дээш үзүүлэлт тайлагнасан байдаг (жишээлбэл FaceNet \cite{schroff2015facenet}, ArcFace \cite{deng2019arcface}). Харин энэ судалгааны 98.35\% нь enrollment open-set gating нөхцөлд, бодит ангийн орчны шуугиантой pipeline дээр гарсан үзүүлэлт гэдгээрээ зорилго, орчин, үнэлгээний протоколоор ялгаатай. Иймээс шууд тоон өрсөлдүүлэлтээс илүү ``бодит орчинд хэрэгжих чадвар''-ын түвшинд харьцуулах нь зөв.

Гэсэн хэдий ч бүрэн хамгаалалттай thesis болгохын тулд дараагийн зайлшгүй ажил нь ground truth-ийг архивлах, олон өдрийн баталгаажуулалт хийх, confidence interval болон paired significance test тооцох, мөн privacy/consent процесстой уялдсан хэрэглээний хувилбарыг боловсруулж дуусгах явдал юм.

\section{Цаашдын ажил}

Цаашдын шатанд дараах чиглэлээр сайжруулалт хийх нь зүйтэй:
\begin{itemize}
    \item manual attendance ground truth-ийг бүрэн архивлаж, FAR/FRR-ийг шууд тооцох;
    \item гэрэлтүүлэгт суурилсан босго тохируулга болон зургийн чанарын тохируулга хийх;
    \item олон анги, олон өдрийн өгөгдлөөр статистикийн баталгаажуулалт хийх;
    \item хэрэглэгчийн интерфейсийн загварыг хэрэглэгчийн тесттэй хамт үнэлэх;
    \item зөвшөөрөл, хадгалалт, устгалын урсгалыг бодит ашиглалтын шаардлагад нийцүүлэх.
\end{itemize}
